\documentclass[11pt]{article}
\usepackage[utf8]{inputenc}
\usepackage[T1]{fontenc}
\usepackage[italian]{babel}
\usepackage{geometry}
\geometry{a4paper, margin=1in}
\usepackage{xcolor}
\usepackage{listings}
\usepackage{hyperref}

\usepackage{tcolorbox}
\tcbuselibrary{breakable}
\begin{document}

\begin{tcolorbox}[colback=blue!5!white,colframe=blue!75!black,title=\textbf{DOMANDA}, breakable]
Un semplice web server
\begin{itemize}
    \item Aprire il file \texttt{serverHTTP.c} in \texttt{Esempi-web/} e analizzarne il contenuto
    \item Compilarlo come nell'esercitazione sull'interfaccia socket ed eseguirlo
    \item Provare ad aprire un secondo terminale nella stessa cartella e a rilanciare lo stesso server. Funziona? Perché?
    \item Aprire il browser preferito e impostare la URL ``\texttt{http://127.0.0.1:8000/}'' Cosa si vede sul browser e sul terminale?
    \item Aprire il browser preferito e impostare la URL ``\texttt{http://localhost:8000/}'' Cosa cambia?
\end{itemize}
\end{tcolorbox}

\begin{tcolorbox}[colback=green!5!white,colframe=green!75!black,title=\textbf{RISPOSTA}, breakable]
\textbf{Soluzione Esercizio 22:} Un semplice web server

\textbf{Operazioni Preliminari e Compilazione}\\
Inizialmente sono stati copiati i file necessari (\texttt{serverHTTP.c}, \texttt{network.c}, \texttt{network.h}) all'interno della cartella dell'esercizio. Analizzando il sorgente di "\texttt{serverHTTP.c}", sono stati corretti alcuni commenti non validi scritti in stile Python (carattere "\#") sostituendoli con quelli corretti in linguaggio C ("\texttt{//}"). Infine, il codice è stato compilato tramite il comando \texttt{gcc}.\\
Il file è un banale server \texttt{TCP} che resta in ascolto sulla porta 8000, accetta la connessione di un client, processa testualmente la richiesta leggendo riga per riga l'header \texttt{HTTP}, invia in risposta una pagina HTML di base e infine chiude la socket.

\textbf{Rilancio del Server su un secondo terminale}\\
Se si prova a lanciare una seconda istanza del server in un altro terminale, mentre il primo è ancora in esecuzione, il programma terminerà subito per un errore. Questo avviene a causa del fallimento della syscall "bind": non è possibile associare la medesima porta e lo stesso indirizzo a due processi in ascolto nello stesso istante.

\textbf{Navigazione tramite l'indirizzo IP (127.0.0.1)}\\
Richiedendo con il browser la URL "\texttt{http://127.0.0.1:8000/}":
\begin{itemize}
    \item \textbf{Nel browser:} Si visualizza il contenuto statico inserito nel server (la scritta "Hello World, this is a very simple HTML document.").
    \item \textbf{Nel terminale:} Il programma stampa la richiesta \texttt{HTTP} pura inviata dal browser. Tra le varie intestazioni è presente la prima riga "\texttt{GET / HTTP/1.1}" e in particolare l'intestazione "\texttt{Host: 127.0.0.1:8000}".
\end{itemize}

\textbf{Navigazione tramite Hostname (localhost)}\\
Richiedendo con il browser la URL "\texttt{http://localhost:8000/}":
\begin{itemize}
    \item \textbf{Nel browser:} Non cambia assolutamente nulla, in quanto l'hostname "localhost" viene risolto localmente nell'IP di loopback 127.0.0.1. Il server riceverà la connessione e invierà in risposta la stessa identica pagina HTML.
    \item \textbf{Nel terminale:} La richiesta catturata cambierà l'intestazione Host, che rifletterà la stringa usata dal client, diventando "\texttt{Host: localhost:8000}". Questo dimostra come l'hostname testuale digitato dall'utente venga inserito dal browser e trasmesso all'interno delle intestazioni \texttt{HTTP}.
\end{itemize}

\textbf{Comandi Utili per il Test e l'Analisi}

\begin{lstlisting}[language=bash, basicstyle=\ttfamily\small]
# 1. Compilazione
gcc serverHTTP.c network.c -o serverHTTP

# 2. Esecuzione del Web Server
./serverHTTP

# 3. Test tramite CLI (utilizzando curl per visualizzare gli header scambiati)
curl -v http://127.0.0.1:8000/
curl -v http://localhost:8000/

# 4. Registrazione del Traffico (cattura su interfaccia locale 'lo' sulla porta 8000)
# Tramite riga di comando (salvataggio su file .pcap):
sudo tshark -i lo -f "tcp port 8000" -w web_traffic.pcap

# 5. Analisi del Traffico Registrato
# --- Tramite Riga di Comando (CLI) ---
# 1. Leggere tutti i pacchetti registrati:
tshark -r web_traffic.pcap

# 2. Filtrare per mostrare solo i messaggi a livello applicativo HTTP:
tshark -r web_traffic.pcap -Y "http"

# 3. Leggere il contenuto completo delle intestazioni (Dettagli pacchetto estesi):
tshark -r web_traffic.pcap -Y "http" -V

# 4. Seguire l'intera conversazione scambiata:
tshark -r web_traffic.pcap -z follow,tcp,ascii,0

# --- Tramite Interfaccia Grafica (GUI) ---
# Lanciare Wireshark normalmente:
sudo wireshark
# Oppure per aprire direttamente il file registrato per l'analisi visiva:
wireshark web_traffic.pcap
\end{lstlisting}
\end{tcolorbox}

\end{document}
